% Elemento de volume em coordenadas cilíndricas
% Feito para o ÉguaTeX (eguatex.com.br). Uso livre (CC0 1.0).
%
% Compile e pronto: sai só a figura, em PDF.
% Para pôr no seu trabalho, use o PDF com \includegraphics ou copie
% o bloco tikzpicture para dentro de uma figure.
%
% Ângulo phi medido no plano xy, a partir do eixo x.
\documentclass[tikz,border=4pt]{standalone}
\usepackage[T1]{fontenc}
\usepackage{lmodern}
\usetikzlibrary{patterns.meta}

% Para usar (cil cs:r=..., p=..., z=...) no seu documento, copie também estas linhas.
\tikzset{cil/r/.store in=\cilr, cil/p/.store in=\cilp, cil/z/.store in=\cilz}
\tikzdeclarecoordinatesystem{cil}{%
  \tikzset{cil/.cd,#1}%
  \pgfpointxyz{\cilr*cos(\cilp)}{\cilr*sin(\cilp)}{\cilz}}

\begin{document}
\begin{tikzpicture}[
  x={(-0.4435cm,-0.3663cm)}, y={(0.9768cm,-0.1715cm)}, z={(0cm,1cm)},
  line cap=round, line join=round,
  eixo/.style={line width=0.75pt},
  aresta/.style={line width=0.8pt},
  guia/.style={line width=0.55pt, black!27}]

% Mude estes números e o desenho e os rótulos acompanham (ângulos em graus).
\def\ra{4.05}  \def\rb{5.4}  % rho e rho + drho
\def\pa{47.75} \def\pb{65}   % phi e phi + dphi
\def\za{2.49}  \def\zb{3.72} % z e z + dz

\draw[eixo] (0,0,0) -- (5.06,0,0) node[below left] {$x$};
\draw[eixo] (0,0,0) -- (0,5.6,0)  node[right]      {$y$};
\draw[eixo] (0,0,0) -- (0,0,4.97) node[above]      {$z$};

\fill[pattern={Lines[distance=3.6pt, line width=0.5pt]}, pattern color=black!25]
  plot[domain=\pa:\pb] (cil cs:r=\ra,p=\x,z=0)
  -- plot[domain=\pb:\pa] (cil cs:r=\rb,p=\x,z=0) -- cycle;
\draw[guia] plot[domain=\pa:\pb] (cil cs:r=\ra,p=\x,z=0);
\draw[guia] plot[domain=\pa:\pb] (cil cs:r=\rb,p=\x,z=0);
\foreach \p in {\pa,\pb} {
  \draw[guia] (0,0,0) -- (cil cs:r=\rb,p=\p,z=0);
  \draw[guia] (0,0,\zb) -- (cil cs:r=\ra,p=\p,z=\zb);
}

\draw[guia] (cil cs:r=\ra,p=\pa,z=0) -- (cil cs:r=\ra,p=\pa,z=\za);
\draw[guia] (cil cs:r=\rb,p=\pa,z=0) -- (cil cs:r=\rb,p=\pa,z=\za);
\draw[guia] (cil cs:r=\ra,p=\pb,z=0) -- (cil cs:r=\ra,p=\pb,z=\za);
\draw[aresta] (cil cs:r=\rb,p=\pb,z=0) -- (cil cs:r=\rb,p=\pb,z=\za);

\filldraw[aresta, fill=black!5]
  (cil cs:r=\ra,p=\pa,z=\za) -- (cil cs:r=\rb,p=\pa,z=\za)
  -- (cil cs:r=\rb,p=\pa,z=\zb) -- (cil cs:r=\ra,p=\pa,z=\zb) -- cycle;
\filldraw[aresta, fill=black!9]
  plot[domain=\pa:\pb] (cil cs:r=\ra,p=\x,z=\zb)
  -- plot[domain=\pb:\pa] (cil cs:r=\rb,p=\x,z=\zb) -- cycle;
\filldraw[aresta, fill=black!14]
  plot[domain=\pa:\pb] (cil cs:r=\rb,p=\x,z=\za)
  -- plot[domain=\pb:\pa] (cil cs:r=\rb,p=\x,z=\zb) -- cycle;
\draw[guia] plot[domain=\pa:\pb] (cil cs:r=\ra,p=\x,z=\za);
\draw[guia] (0,0,\za) -- (cil cs:r=\rb,p=\pb,z=\za);
\draw[guia] (cil cs:r=\ra,p=\pb,z=\za) -- (cil cs:r=\ra,p=\pb,z=\zb);

\draw[aresta] (0,0,\za) -- (cil cs:r=\ra,p=\pa,z=\za);
\fill (cil cs:r=\ra,p=\pa,z=\za) circle[radius=1.9pt];
\draw[eixo] plot[domain=0:\pa]   (cil cs:r=1.65,p=\x,z=0);
\draw[eixo] plot[domain=\pa:\pb] (cil cs:r=1.95,p=\x,z=0);

\path (0,0,\za) -- (cil cs:r=\ra,p=\pa,z=\za) node[midway, above] {$\rho$};
\path (cil cs:r=\ra,p=\pa,z=\za) -- (cil cs:r=\rb,p=\pa,z=\za)
  node[midway, below left] {$\mathrm{d}\rho$};
\node at (cil cs:r={(\ra+\rb)/2},p={(\pa+\pb)/2},z={(\za+\zb)/2}) {$\mathrm{d}V$};
\path (cil cs:r=\rb,p=\pb,z=\za) -- (cil cs:r=\rb,p=\pb,z=\zb) node[midway, right] {$\mathrm{d}z$};
\path (cil cs:r=\rb,p=\pb,z=0)   -- (cil cs:r=\rb,p=\pb,z=\za) node[midway, right] {$z$};
\node[below] at (cil cs:r=\rb,p={(\pa+\pb)/2},z=\za) {$\rho\,\mathrm{d}\phi$};
\node at (cil cs:r=1.2,p={\pa/2},z=0) {$\phi$};
\node at (cil cs:r=2.45,p={(\pa+\pb)/2},z=0) {$\mathrm{d}\phi$};

\node[align=center] at (3.6cm,4.6cm) {\textit{Coordenadas cilíndricas}\\[2pt] $(\rho,\phi,z)$};
\node at (1.7cm,-3.2cm) {$\mathrm{d}V=\rho\,\mathrm{d}\rho\,\mathrm{d}\phi\,\mathrm{d}z$};
\end{tikzpicture}
\end{document}
